%% sections/slide10-acknowledgements.tex — SLIDE 10 — Acknowledgements
%% =============================================================================
%% CONTENT BLUEPRINT — Acknowledgements
%% =============================================================================
%% Time budget: read-only, 15 seconds. Same role as references.
%%
%% WHAT TO PUT:
%%
%%   • Supervisor first (most formal), co-supervisor if applicable.
%%   • Specific person who lent equipment / sample / scientific guidance.
%%     Vague "my lab" thanks are read as boilerplate by the examiner —
%%     name them.
%%   • Funding: which grant, fellowship, or departmental fund supported
%%     the work. \highlight{Optional but recommended}.
%%   • Family / friends — neutral, brief.
%%
%% STYLE:
%%   • One short paragraph, no more than 5 short sentences.
%%   • Use a centered format.
%%   • Names italicised (SUST physics style).
%%
%% CASUAL-ASIDE FORM (optional, but only if you have rapport with the
%% examiner): a small thanks to peers you worked with day-to-day. Helps
%% the room relax after a dense slides 3–7.
%% =============================================================================

\begin{frame}{Acknowledgements}
  \begin{center}
    \vspace{-2pt}
    \itshape
    First and foremost, I thank my supervisor,
    \textbf{\defSupervisor}, for \highlight{[specific contribution:
    guidance / supervision of the experiment / reading this manuscript /
    weekly discussions over the past [$N$] months]}.

    \medskip
    I am equally grateful to \highlight{[specific person: lab technician /
    collaborator / second supervisor]}, who \highlight{[specific
    contribution: provided the SF59 sample / ran DFT baseline / wrote the
    fitting routine]}.

    \medskip
    This work was supported by \highlight{[grant / fellowship / department
    fund]}. Equipment access at \highlight{[facility]} is acknowledged.

    \medskip
    Finally, thanks to my family and peers for their support throughout
    this work.
  \end{center}
\end{frame}
